\section{A quartic obstruction to a benign optimization landscape}
The previous note ruled out a naive local-maximization argument for
the entrywise one-norm. We now test the smoother alternative
\[
 Q(U)=\sum_{ij}u_{ij}^4\qquad(U\in O(n)).
\]
It satisfies $Q(U)\ge1$, with equality exactly when $\sqrt nU$ is a
Hadamard matrix. The arbitrary-exponent framework is classical
\cite{ArbitraryP}; the calculations below are independently derived,
and no novelty claim is made.

\begin{proposition}[An infinite family of strict quartic traps]
Let $A=2J_5/5-I_5$, and let $B=H/\sqrt k$ for any Hadamard matrix
$H$ of order $k$. Then $U=A\otimes B$ is a strict local minimum of
$Q$ on $O(5k)$, with value $29/25$. In fact its Hessian obeys
\[
 \left.\frac{d^2}{dt^2}Q(Ue^{tX})\right|_{t=0}
 \ge\frac4{5k}\|X\|_F^2
 \qquad(X^{\mathsf T}=-X).
\]
For $k=4\cdot2^a$, these are nonglobal local minima at the admissible
orders $20\cdot2^a$.
\end{proposition}
\begin{proof}
For an orthogonal matrix $V$, entrywise differentiation gives
\begin{align*}
 \left.\frac{d}{dt}Q(Ve^{tX})\right|_{0}
 &=4\operatorname{tr}((V^{\circ3})^{\mathsf T}VX),\\
 \left.\frac{d^2}{dt^2}Q(Ve^{tX})\right|_{0}
 &=12\sum_{ij}v_{ij}^2(VX)_{ij}^2
   +4\operatorname{tr}((V^{\circ3})^{\mathsf T}VX^2).
\end{align*}
At the claimed $U$, the symmetric matrix
\[
 M=U^{\mathsf T}U^{\circ3}
  =\left(\frac7{25}I_5-\frac6{125}J_5\right)\otimes\frac1kI_k
\]
has eigenvalues $1/(25k)$ and $7/(25k)$. Thus $U$ is stationary.
Its smallest squared entry is $4/(25k)$. Since
$\|UX\|_F=\|X\|_F$ and $X^2=-X^{\mathsf T}X$,
\[
 Q''(0)\ge\left(\frac{48}{25k}-\frac{28}{25k}\right)\|X\|_F^2
 =\frac4{5k}\|X\|_F^2.
\]
The positive definite tangent Hessian proves strict local minimality.
Direct summation gives $Q(A)=29/25$ and $Q(B)=1$, hence the stated
value. A Hadamard matrix of order $20$ is furnished by the certified
length-five Goethals--Seidel witness; tensoring it with the Sylvester
matrix of order $2^a$ gives a point of value one at the same order.
Consequently these local minima are not global.
\end{proof}

The obstruction is uniform, not an inference from numerical Hessian
eigenvalues. The exact script checks the algebra at orders $20,40,80$.
It shows that changing the one-norm objective to the fourth-power
objective does not by itself remove the false-minimum problem.

\begin{proposition}[Coordinate rotations can miss a descent direction]
At order $28$, let
$U=(2J_7/7-I_7)\otimes(H_4/2)$. This matrix is stationary for $Q$.
Every single Givens rotation has strictly positive second derivative,
but a joint tangent direction has negative second derivative.
\end{proposition}
\begin{proof}
For $A_n=2J_n/n-I_n$, direct substitution in the Hessian formula yields
\[
 \left.\frac{d^2}{dt^2}Q(A_ne^{tX})\right|_0
 =-\frac{4(n-6)}n\|X\|_F^2
  +\frac{8(n^2-16)}{n^3}\|X\one\|^2.
\]
Take $n=7$. A single Givens direction in the tensor matrix has
curvature $34/343$ when its two coordinates belong to distinct
order-seven blocks, and $412/343$ when they belong to the same block.
Both are positive. Set
$X_0=E_{01}+E_{12}-E_{02}$, where
$E_{ij}=e_ie_j^{\mathsf T}-e_je_i^{\mathsf T}$. Then
$X_0\one=0$ and $\|X_0\|_F^2=6$. The joint direction
$X=X_0\otimes I_4$ gives curvature $-24/7$, since the fourth-power
objective of the normalized Hadamard tensor factor is one.
All values are also checked by the exact full tangent Hessian script.
\end{proof}
Thus testing all coordinate-pair second derivatives is not a sufficient
second-order test even at an admissible order. This matters when
importing coordinate descent or block-rotation methods from numerical
optimization. The first proposition supplies genuine strict minima;
the second supplies a different failure, namely a saddle hidden from
every individual coordinate rotation.
